\section*{Question 8}
If ability is omitted, the regression becomes
\[
\mathrm{wage}_i=\beta_0+\beta_1\mathrm{education}_i+v_i,
\qquad v_i=\beta_2\mathrm{ability}_i+u_i.
\]
Even if the full model satisfies
$E[u_i\mid\mathrm{education}_i,\mathrm{ability}_i]=0$, the short regression's
zero conditional mean assumption, $E[v_i\mid\mathrm{education}_i]=0$,
can fail: ability is part of the error and is correlated with education.
Allowing the intercept to absorb a constant error mean does not remove
this dependence on education.

Under the full model's exogeneity assumption, the omitted-variable bias
formula is
\[
\operatorname{plim}\tilde{\beta}_1-\beta_1
=\beta_2\frac{\operatorname{Cov}(\mathrm{education}_i,\mathrm{ability}_i)}
{\operatorname{Var}(\mathrm{education}_i)}
=\beta_2\delta_1,
\]
where $\delta_1$ is the population slope from regressing ability on education.
If ability raises wages ($\beta_2>0$) and is positively associated with
education ($\delta_1>0$), the bias is positive. Omitting ability then
attributes some of its wage association to education, so the education
coefficient is expected to be larger than in the model controlling for
ability. The sample estimates are consistent with this reasoning:
4.760 without ability versus 3.690 with ability. This conclusion about bias
relies on the full model's assumptions; including ability alone does not
guarantee that all confounding has been removed.
